\documentclass[11pt]{article}

\usepackage{fullpage}
\usepackage{fourier}
\usepackage{graphicx}
\usepackage{xspace}
\usepackage{epigraph}
\usepackage{listings}
\usepackage{xcolor}
\usepackage{url}
\usepackage{soul}
\usepackage{multirow}
\usepackage{clrscode3e}
\usepackage{hyperref}
\usepackage{wrapfig}
\usepackage{amsmath}

\hypersetup{
    colorlinks=false
}

\usepackage{fancyhdr}
\pagestyle{fancy}
\fancyhf{}

\fancypagestyle{plain}{%
  \fancyhf{}
  \renewcommand{\headrulewidth}{0pt}
  \renewcommand{\footrulewidth}{0pt}
  \lfoot{\textcopyright{} 2021 Darrell Long}
  \rfoot{\thepage}
}

\pagestyle{plain}

\definecolor{codegreen}{rgb}{0,0.5,0}
\definecolor{codegray}{rgb}{0.5,0.5,0.5}
\definecolor{codepurple}{rgb}{0.58,0,0.82}

\lstloadlanguages{C,make,python,fortran}

\lstdefinestyle{c99}{
    morekeywords={bool, uint8_t, uint16_t, uint32_t, uint64_t, int8_t, int16_t, int32_t, int64_t},
    commentstyle=\color{codegreen},
    keywordstyle=\color{magenta},
    numberstyle=\tiny\color{codegray},
    identifierstyle=\color{blue},
    stringstyle=\color{codepurple},
    basicstyle=\ttfamily,
    breakatwhitespace=false,
    breaklines=true,
    captionpos=b,
    keepspaces=true,
    numbers=left,
    numbersep=5pt,
    showspaces=false,
    showstringspaces=false,
    showtabs=false,
    tabsize=4
}


\title{Assignment 5 \\ Hamming Codes}
\author{Prof. Darrell Long \\ CSE 13S -- Spring 2021}
\date{
  First \texttt{DESIGN.pdf} draft due: May 6$^\text{th}$ at 11:59\,pm PST \\
  Assignment due: May 9$^\text{th}$ at 11:59\,pm PST
}

\begin{document}\maketitle

\section{Introduction}

As we know, the world is far from perfect. In the communications domain,
this imperfection is called \emph{noise}. Noise (unwanted random
disturbances) makes it difficult to have a reliable signal. Thus,
transferring data through a noisy communication channel is prone to
errors. Noisy channels are omnipresent. They are present in mobile
phone networks and even the wires in a circuit. To counteract noisy
interference, we add extra information to our data. This extra
information allows us to perform error checking, and request that the
sender retransmit any data that was incorrect. We can also add extra
information to not only detect errors but also correct them. This
technique is called \emph{forward error correction} (FEC). CDs, DVDs,
and even hard drives use FEC to account for scratches or bad sectors. In
fact, most of our digital world such as Netflix would not be possible
without FEC.

\section{Hamming Codes}

\epigraphwidth=0.75\textwidth
\epigraph{
  Hamming: \emph{If you come to the Navy School, I will teach you how to be a great scientist.}\\
  Long: \emph{And if I go to Santa Cruz, will you still teach me?}\\
  Hamming: \emph{No.}
}{Lunch with Richard Hamming}

\begin{wrapfigure}{r}{0.15\textwidth}
\includegraphics[width=0.15\textwidth]{images/hamming.jpeg}
\centerline{\small Richard Hamming}
\end{wrapfigure}

Richard Wesley Hamming was an American applied mathematician whose work
had profound impact on computer engineering and telecommunications. His
contributions include the Hamming code (which makes use of a Hamming
matrix), the Hamming window, Hamming numbers, sphere-packing (or Hamming
bound), and the Hamming distance. Hamming served as president of the
Association for Computing Machinery from 1958 to 1960.

He used to joke that he was the \emph{anti-Huffman}: David Huffman, the
inventor of Huffman codes and a professor here at Santa Cruz, was a
friend. Hamming added redundancy for reliability, while Huffman focused
on removing it for efficiency.

In later life, Hamming became interested in teaching. From 1960 and
1976, before he left Bell Laboratories, he held visiting positions at
Stanford University, Stevens Institute of Technology, the City College
of New York, the University of California at Irvine and Princeton
University. After retiring from the Bell Laboratories in 1976, Hamming
took a position at the Naval Postgraduate School in Monterey,
California, where he worked as an adjunct professor and senior lecturer
in computer science, and devoted himself to teaching and writing books.

He spent significant effort in trying to recruit a young Dr.\xspace
Long, telling him that ``If you come to the Navy School, I will teach
you how to be a great scientist.'' Dr.\xspace Long replied, ``If I go to
Santa Cruz, will you still teach me?'' To which Hamming replied, simply,
``No.'' Sadly, it was an opportunity lost.

\subsection{Bits, Nibbles, and Bytes}

Around 1948, John Wilder Tukey, an American mathematician, coined the
word \emph{bit} to replace the mouthful that was \emph{binary digit}. A
bit is the basic unit of information for digital systems. It represents
a logical state with only two possible values, either a 1 or 0, on or
off, true or false. But one bit is rather limiting. As a result,
multiple bits were packed together to make up a \emph{nibble} (4 bits
with $2^4$ states) or a \emph{byte} (8 bits with $2^8$ states).

\subsection{Overview}

Hamming codes are a linear error-correcting code invented by Richard W.
Hamming \footnote{R. W. Hamming, ``Error detecting and error correcting
codes,'' \textit{The Bell System Technical Journal}, April 1950.} to
correct errors caused by punched card readers.  But, we will be using
them to correct errors caused by random bit-flips (noise). He introduced
the Hamming(7,4) code that encodes 4 bits of data $D$ in 7 bits by
adding 3 redundant or parity $P$ bits. An explanation of parity bits
will be given through an example. This code can detect and correct
\emph{one} error. With 7 possible errors, 3 parity bits can identify
which bit is incorrect ($2^3 - 1 = 7$).

The Hamming(7,4) code is shown in table \ref{ham7,4} where the Hamming
code's \emph{least significant bit} (LSB) is at index $001_2$ and the
\emph{most significant bit} (MSB) is in position $111_2$. While normally
in practice we begin counting at 0, the importance of starting the count
at 1 and in binary will be evident in the next paragraphs.

\begin{table}[h]
  \centering
  \begin{tabular}{ |c|c|c|c|c|c|c|c| }
     \hline
     Index & $111_2$ & $110_2$ & $101_2$ & $100_2$ & $011_2$ & $010_2$ & $001_2$ \\
     \hline \hline
     Hamming code & $D_3$ & $D_2$ & $D_1$ & $P_2$ & $D_0$ & $P_1$ & $P_0$ \\
     \hline
  \end{tabular}
  \caption{Hamming(7,4) code with data bits $D$ and parity bits $P$.}
  \label{ham7,4}
\end{table}


You might notice that each parity bit's index has only one bit set (a
power of 2). $P_0$'s index ($001_2$) has the $0^\text{th}$ bit set, and you
might notice that the index for $D_0$, $D_1$, and $D_3$ also have the
$0^\text{th}$ bit set. Thus, $P_0$ is calculated over $D_0$, $D_1$, and
$D_3$, \textit{i.e.}, $P_0$ is set if the data bits have an odd number
of 1s. This is an \emph{even} parity scheme. With an odd parity scheme,
the parity bit is set if by setting it there are an odd number of 1s. We
could use an odd parity scheme, but for this assignment, we will be
using an even parity.

We can calculate $P_1$ and $P_2$ in the same way where $P_i$ is
calculated over the data bits whose index has the $i^\text{th}$ bit set. The
formulas for the three parity bits $P_i$ are shown below.
\begin{equation*}
  \begin{split}
    P_0 &= D_0 \oplus D_1 \oplus D_3 \\
    P_1 &= D_0 \oplus D_2 \oplus D_3 \\
    P_2 &= D_1 \oplus D_2 \oplus D_3
  \end{split}
\end{equation*}

\noindent As a result, each data bit has at least two parity bits that
will help recover its value should an error occur. The overlap of parity
bits also keeps them in check (parity bits can be erroneously flipped
too). For example, the Hamming code for $0001_2$ is shown in table
\ref{ham7,4-ex}.  From here on, the index will follow convention and
start at 0 and in decimal.

\begin{table}[h]
  \centering
  \begin{tabular}{ |c|c|c|c|c|c|c|c| }
    \hline
    Index & 6 & 5 & 4 & 3 & 2 & 1 & 0 \\
    \hline \hline
    Label & $D_3$ & $D_2$ & $D_1$ & $P_2$ & $D_0$ & $P_1$ & $P_0$ \\
    \hline
    Hamming code & $0$ & $0$ & $0$ & $0$ & $1$ & $1$ & $1$ \\
    \hline
  \end{tabular}
  \caption{Hamming(7,4) code for $0001_2$.}
  \label{ham7,4-ex}
\end{table}

\noindent The values of parity bits $P_0$, $P_1$, and $P_2$ are
calculated as follows:
\begin{equation*}
  \begin{split}
    P_0 &= D_0 \oplus D_1 \oplus D_3 = 1 \oplus 0 \oplus 0 = 1 \\
    P_1 &= D_0 \oplus D_2 \oplus D_3 = 1 \oplus 0 \oplus 0 = 1 \\
    P_2 &= D_1 \oplus D_2 \oplus D_3 = 0 \oplus 0 \oplus 0 = 0
  \end{split}
\end{equation*}

We currently have a \emph{non-systematic} code, a code where the parity
bits are interspersed throughout the code. While this is fine for a
hardware implementation, it is tedious to do in software.  Instead, we
will be using a \emph{systematic} code, a code where the parity bits are
placed after the data bits. We will also be extending the Hamming(7,4)
code by adding one more parity bit to make this a Hamming(8,4) code. The
extra parity bit, $P_3$, is calculated over $P_0$--$P_2$ and
$D_0$--$D_3$. If $P_3$ for a received code is incorrect then we know an
error has occurred. Either $P_3$ is incorrect or one of the the bits in
the code is incorrect. This extension has two additional benefits.
First, each code is a byte rather than seven bits. Second, we can now
detect two errors but only correct one. The new Hamming code is shown
below.

\begin{table}[h]
  \centering
  \begin{tabular}{ |c|c|c|c|c|c|c|c|c| }
    \hline
    Index & 7 & 6 & 5 & 4 & 3 & 2 & 1 & 0 \\
    \hline
    \hline
    Hamming code & $P_3$ & $P_2$ & $P_1$ & $P_0$ & $D_3$ & $D_2$ & $D_1$ & $D_0$ \\
    \hline
  \end{tabular}
  \caption{Hamming(8,4) systematic code, an extension of the Hamming(7,4) code.}
  \label{ham8,4}
\end{table}

\noindent The values of parity bits $P_0$, $P_1$, $P_2$, and $P_3$ are
calculated as follows:
\begin{equation*}
\begin{split}
    P_0 &= D_1 \oplus D_2 \oplus D_3 \\
    P_1 &= D_0 \oplus D_2 \oplus D_3 \\
    P_2 &= D_0 \oplus D_1 \oplus D_3 \\
    P_3 &= D_0 \oplus D_1 \oplus D_2 \oplus D_3 \oplus P_0 \oplus P_1
    \oplus P_2 = D_0 \oplus D_1 \oplus D_2 \\
\end{split}
\end{equation*}

\subsection{Encoding}

One approach to generating a Hamming code for a message is to calculate
the parity bits one-by-one and appending them to the end of the message.
Instead, we can use a generator matrix, $\boldsymbol{G}$.  Given a
message, $\Vec{m}$, of four bits (a nibble) we can generate its Hamming
code, $\Vec{c}$, by vector-matrix multiplication $\Vec{c} =
\Vec{m}\boldsymbol{G}$ where the resulting code is eight bits in size (a
byte). $\boldsymbol{G}$ is defined as:

\begin{equation*}
  \boldsymbol{G} =
  \begin{pmatrix}
    1 & 0 & 0 & 0 & 0 & 1 & 1 & 1 \\
    0 & 1 & 0 & 0 & 1 & 0 & 1 & 1 \\
    0 & 0 & 1 & 0 & 1 & 1 & 0 & 1 \\
    0 & 0 & 0 & 1 & 1 & 1 & 1 & 0 \\
  \end{pmatrix}.
\end{equation*}

\noindent Vector-matrix multiplication for a vector $\Vec{y}$ of size
$n$ and matrix $\boldsymbol{A}$ of size $n \times n$ where $\Vec{y} =
\Vec{x}\boldsymbol{A}$ is defined as:
\begin{equation*}
  \begin{split}
    y_{i} = \sum\limits_{k=1}^n x_k \cdot A_{i,k}
  \end{split}.
\end{equation*}

Those of you with exposure to linear algebra (helpful but not required)
will notice the left half of $\boldsymbol{G}$ is the identity matrix
\textit{I} (1s along the diagonal). This ensures the first four bits of
the Hamming code are the data bits while the right half is used to
calculate the parity of message (notice the 0s along the diagonal).
Generating the Hamming Code for $\Vec{m} = \begin{pmatrix} 0 & 0 & 1 & 1
\end{pmatrix}$ ($1100_2$ in binary) is shown below.
\textcolor{red}{Note: binary is read from right to left with the LSB in
the rightmost position and the MSB in the leftmost position. Vectors
(arrays) are read from left to right.}

\begin{align*}
    \Vec{c} &=
      \bordermatrix{~ & 0 & 1 & 2 & 3 \cr
        ~ & 0 & 0 & 1 & 1 \cr}
    \begin{pmatrix}
    1 & 0 & 0 & 0 & 0 & 1 & 1 & 1 \\
    0 & 1 & 0 & 0 & 1 & 0 & 1 & 1 \\
    0 & 0 & 1 & 0 & 1 & 1 & 0 & 1 \\
    0 & 0 & 0 & 1 & 1 & 1 & 1 & 0
  \end{pmatrix} \pmod{2} \\
   &= \begin{pmatrix}0 & 0 & 1 & 1 & 2 & 2 & 1 & 1\end{pmatrix} \pmod{2} \\
   &=\begin{pmatrix}0 & 0 & 1 & 1 & 0 & 0 & 1 & 1\end{pmatrix}
\end{align*}

\noindent In binary, $\Vec{c}$ is equivalent to $1100$ $1100_2$.

\textcolor{red}{Note: Addition and multiplication is $\mathrm{mod}\ 2$ in
binary.} You might remember from CSE 12 that $\bigoplus$ (exclusive-or)
is the summing function and can be used in lieu of addition, and
$\wedge$ (logical and) can substitute multiplication. These operations
have the added benefit of operating in mod 2.

\begin{equation*}
    \begin{split}
        a \oplus b &= a + b \pmod{2} \\
        a \wedge b &= a \times b \pmod{2}
    \end{split}
\end{equation*}

\subsection{Decoding}

The process for decoding a Hamming code is similar to encoding a message
as it uses a parity-checker matrix, $\boldsymbol{H}$, to identify any
errors and recover the original message if an error occurred.  To decode
a message, the code is multiplied by the transpose of the parity-checker
matrix, $\Vec{e} = \Vec{c}\boldsymbol{H^\intercal}$ where $\Vec{e}$ is
the \emph{error syndrome}, $\Vec{c}$ is the Hamming code, and
$\boldsymbol{H}$ is defined as

\begin{equation*}
  \boldsymbol{H} =
  \begin{pmatrix}
    0 & 1 & 1 & 1 & 1 & 0 & 0 & 0 \\
    1 & 0 & 1 & 1 & 0 & 1 & 0 & 0 \\
    1 & 1 & 0 & 1 & 0 & 0 & 1 & 0 \\
    1 & 1 & 1 & 0 & 0 & 0 & 0 & 1
  \end{pmatrix}.
\end{equation*}

However, $\Vec{e} \neq \Vec{m}$.  Instead, $\Vec{e}$, the \emph{error
syndrome}, is a pattern of bits that identifies the error if there is
one. For example, if $\Vec{e} = [0, 1, 1, 1]$, matching
$\boldsymbol{H}^\intercal$'s first row, then we know the error lies in
the $0^\text{th}$ bit and flipping the $0^\text{th}$ bit in $\Vec{c}$
will correct the error (remember the first four bits in the Hamming code
are the data bits). If $\Vec{e} = \boldsymbol{0}$ then our message does
not contain an error. But, if the error syndrome is non-zero and the
vector is not one of $\boldsymbol{H}^\intercal$'s rows then we cannot
correct the error since more than one bit has been flipped.

For example, to decode $\Vec{c}$ calculated earlier, we can do the following:

\begin{align*}
  \Vec{e}&=\Vec{c}\boldsymbol{H^{\intercal}} =
  \bordermatrix{
    ~ & 0 & 1 & 2 & 3 & 4 & 5 & 6 & 7\cr
    ~ & 0 & 0 & 1 & 1 & 0 & 0 & 1 & 1\cr
  }
  \begin{pmatrix}
     0 & 1 & 1 & 1 \\
     1 & 0 & 1 & 1 \\
     1 & 1 & 0 & 1 \\
     1 & 1 & 1 & 0 \\
     1 & 0 & 0 & 0 \\
     0 & 1 & 0 & 0 \\
     0 & 0 & 1 & 0 \\
     0 & 0 & 0 & 1
  \end{pmatrix} \pmod{2} \\
  &= \begin{pmatrix}0 & 0 & 0 & 0 \end{pmatrix}.
\end{align*}

\noindent Since $\Vec{e} = \begin{pmatrix}0 & 0 & 0 & 0 \end{pmatrix}$,
we know our message arrived with no errors. Since this is a systematic
code, $\Vec{m} = \begin{pmatrix}0 & 0 & 1 & 1 \end{pmatrix}$ (the first
four elements in $\Vec{c}$).

If the second bit had been flipped due to noise and we received $\Vec{c}
= \begin{pmatrix}0 & 1 & 1 & 1 & 0 & 0 & 1 & 1\end{pmatrix}$ instead, we
can calculate $\Vec{e}$ to identify the flipped bit.

\begin{align*}
    \Vec{e} &=\Vec{c}\boldsymbol{H^{\intercal}} =\begin{pmatrix}0 & 1 & 1 & 1 & 0 & 0 & 1 & 1\end{pmatrix}
    \begin{pmatrix}
         0 & 1 & 1 & 1 \\
         1 & 0 & 1 & 1 \\
         1 & 1 & 0 & 1 \\
         1 & 1 & 1 & 0 \\
         1 & 0 & 0 & 0 \\
         0 & 1 & 0 & 0 \\
         0 & 0 & 1 & 0 \\
         0 & 0 & 0 & 1
    \end{pmatrix} \pmod{2} \\
    &= \begin{pmatrix}1 & 0 & 1 & 1 \end{pmatrix}.
\end{align*}

\noindent Since $\Vec{e} = \begin{pmatrix}1 & 0 & 1 & 1 \end{pmatrix}$,
which is $\boldsymbol{H}^\intercal$'s second row, we know the second
element in $\Vec{c}$ was erroneously flipped. Flipping the value of the
second element gives us $\Vec{c} = \begin{pmatrix}0 & 0 & 1 & 1 & 0 & 0
& 1 & 1\end{pmatrix}$.  As a result, $\Vec{m} = \begin{pmatrix}0 & 0 & 1
& 1\end{pmatrix}$ or $1100_2$, the original message. One possible method
to determine which bit to flip is to compare the error syndrome with
each row in $\boldsymbol{H}^\intercal$, but section~\ref{ham_dec} will go over
an optimal approach involving lookup tables.

\section{Bit Vector}

A \textit{\textbf{bit vector}} is an ADT that represents a one
dimensional array of bits, the bits in which are used to denote if
something is true or false (1 or 0). This is an efficient ADT since, in
order to represent the truth or falsity of a bit vector of $n$ items, we
can use $\lfloor{n / 8}\rfloor + 1$ \texttt{uint8\_t}s instead of $n$,
and being able to access $8$ indices with a single integer access is
extremely cost efficient. Since we cannot directly access a bit, we must
use bitwise operations to get, set, and clear a bit within a byte.

\begin{codelisting}{}
struct BitVector {
    uint32_t length; // Length in bits.
    uint8_t *vector; // Array of bytes.
};
\end{codelisting}

\subsection{\texttt{BitVector *bv\_create(uint32\_t length)}}

The constructor for a bit vector. In the event that sufficient memory
cannot be allocated, the function must return \texttt{NULL}. Else, it
must return a \texttt{(BitVector *)}, or a pointer to an allocated
\texttt{BitVector}. Each bit of the bit vector should be initialized to
0.

\subsection{\texttt{void bv\_delete(BitVector **v)}}

The destructor for a bit vector. Remember to set the pointer to
\texttt{NULL} after the memory associated with the bit vector is freed.

\subsection{\texttt{uint32\_t bv\_length(BitVector *v)}}

Returns the length of a bit vector.

\subsection{\texttt{void bv\_set\_bit(BitVector *v, uint32\_t i)}}

Sets the $i^\text{th}$ bit in a bit vector. To set a bit in a bit
vector, it is necessary to first locate the byte in which the
$i^\text{th}$ resides. The location of the byte is calculated as
\texttt{i / 8}. The position of the bit in that byte is calculated as
\texttt{i \% 8}. Setting a bit in a bit vector should not modify any of
the other bits.

\subsection{\texttt{void bv\_clr\_bit(BitVector *v, uint32\_t i)}}

Clears the $i^\text{th}$ bit in the bit vector.

\subsection{\texttt{uint8\_t bv\_get\_bit(BitVector *v, uint32\_t i)}}

Returns the $i^\text{th}$ bit in the bit vector.

\subsection{\texttt{void bv\_xor\_bit(BitVector *v, uint32\_t i, uint8\_t bit)}}

\texttt{XOR}s the $i^\text{th}$ bit in the bit vector with the value of
the specified \texttt{bit}.

\subsection{\texttt{void bv\_print(BitVector *v)}}

A debug function to print a bit vector. You will want to do this first
in order to verify the correctness of your bit vector implementation.

\section{Bit Matrix}

In this assignment, we will define a bit matrix ADT, which will be an
abstraction over a bit vector. That is, given an $m \times n$ bit matrix
$M$, $0 \leq r < m$, and $0 \leq c < n$, then $M_{r, c}$ corresponds
with index $\texttt{r * n + c}$ in the underlying bit vector. Using this
bit matrix abstraction will help when performing matrix multiplication
$\mathrm{mod}\ 2$ when encoding and decoding Hamming codes.

\begin{codelisting}{}
struct BitMatrix {
    uint32_t rows;
    uint32_t cols;
    BitVector *vector;
};
\end{codelisting}

\subsection{\texttt{BitMatrix *bm\_create(uint32\_t rows, uint32\_t cols)}}

The constructor for a bit matrix. In the event that sufficient memory
cannot be allocated, the function must return \texttt{NULL}. Else, it
must return a \texttt{(BitMatrix *)}, or a pointer to an allocated
\texttt{BitMatrix}. The number of bits in the bit matrix is calculated
as \texttt{row} $\times$ \texttt{cols}. Each bit should be initialized
to 0.

\subsection{\texttt{void bm\_delete(BitMatrix **m)}}

The destructor for a bit matrix. Remember to set the pointer to
\texttt{NULL} after the memory associated with the bit matrix is freed.

\subsection{\texttt{uint32\_t bm\_rows(BitMatrix *m)}}

Returns the number of rows in the bit matrix.

\subsection{\texttt{uint32\_t bm\_cols(BitMatrix *m)}}

Returns the number of columns in the bit matrix.

\subsection{\texttt{void bm\_set\_bit(BitMatrix *m, uint32\_t r, uint32\_t c)}}

Sets the bit at row \texttt{r} and column \texttt{c} in the bit matrix.
Remember, given an $m \times n$ bit matrix $M$, $0 \leq r < m$, and $0
\leq c < n$, then $M_{r, c}$ corresponds with index \texttt{r * n + c}
in the underlying bit vector. Setting a bit should not modify any of the
other bits.

\subsection{\texttt{void bm\_clr\_bit(BitMatrix *m, uint32\_t r,
uint32\_t c)}}

Clears the bit at row \texttt{r} and column \texttt{c}. Clearing a bit
should not modify any of the other bits.

\subsection{\texttt{uint8\_t bm\_get\_bit(BitMatrix *m, uint32\_t r, uint32\_t c)}}

Gets the bit at row \texttt{r} and column \texttt{c}. Getting a bit
should not modify any of the other bits.

\subsection{\texttt{BitMatrix *bm\_from\_data(uint8\_t byte, uint32\_t length)}}

We will occasionally need to transform the first \texttt{length} number
of bits in a byte into a bit matrix. The returned bit matrix should have
the dimensions $1 \times \texttt{length}$. The value of \texttt{length}
should never be greater than 8, since we will only need to transform
codes (8 bits) or nibbles of data (4 bits) into bit matrices.

\subsection{\texttt{uint8\_t bm\_to\_data(BitMatrix *m)}}

Extracts the first 8 bits of a bit matrix, returning those bits as a
\texttt{uint8\_t}. Hint: since there are 8 bits in a byte, what would be
the most efficient way of getting these bits?

\subsection{\texttt{BitMatrix *bm\_multiply(BitMatrix *A, BitMatrix *B)}}

Performs a true matrix multiply, multiplying bit matrix \texttt{A} and
bit matrix \texttt{B} $\mathrm{mod}\ 2$. Returns a new bit matrix that
holds the result of the multiplication. Remember, given an $m \times n$
matrix $A$  and an $n \times p$ matrix $B$, the $m \times p$ matrix
product $C$ is defined as follows for $1 \le i \le m$ and $1 \le j
\le p$:

\[
  C_{i,j} = \sum_{k = 1}^{n} A_{i,k} \times B_{k,j}.
\]

\subsection{\texttt{void bm\_print(BitMatrix *m)}}

A debug function to print a bit matrix. You will want to do this first
in order to verify the correctness of your bit matrix implementation.

\section{Hamming Code Module}

This module will implement the Hamming(8,4) code described in \S 2.2.
The API is defined in \texttt{hamming.h}. As with all provided header
files, \textcolor{red}{you may not modify \texttt{hamming.h}}. The
implementation of the module should be defined in \texttt{hamming.c}.

\subsection{\texttt{HAM\_STATUS}}

The module will contain enumerated status codes that reflect the
different possible cases that can occur when decoding. \texttt{HAM\_OK}
will indicate that no corrections were needed when decoding.
\texttt{HAM\_CORRECT} will indicate that an error was detected and
corrected when decoding. Lastly, \texttt{HAM\_ERR} will indicate that
more than error was detected and a code could not be properly decoded.
These status codes are provided in \texttt{hamming.h}.

\begin{codelisting}{}
typedef enum HAM_STATUS {
    HAM_OK      = -3,  // No error detected.
    HAM_ERR     = -2,  // Uncorrectable.
    HAM_CORRECT = -1   // Detected error and corrected.
} HAM_STATUS;
\end{codelisting}

%\subsubsection{\texttt{ham\_rc ham\_init(void)}}
%This should be called before any function in the module as it creates and initializes
%the $\boldsymbol{G}$ and $\boldsymbol{H}$ Bit Matrix ADTs. If creating either of them
%fails, it must return \texttt{HAM\_ERR} or \texttt{HAM\_OK} otherwise.
%
%\subsubsection{\texttt{void ham\_rc ham\_destroy(void)}}
%To prevent any memory leaks, the Hamming module must free any memory it allocated.
%To check for memory leaks, you must use \texttt{valgrind}.
%\textcolor{red}{Your program must pass valgrind cleanly.}

\subsection{\texttt{uint8\_t ham\_encode(BitMatrix *G, uint8\_t msg)}}

Generates a Hamming code given a nibble of data stored in the lower of
nibble of \texttt{msg} and the generator matrix $\boldsymbol{G}$.
Returns the generated Hamming code, which is stored as a byte, or a
\texttt{uint8\_t}.

\subsection{\texttt{HAM\_STATUS ham\_decode(BitMatrix *Ht, uint8\_t
code, uint8\_t *msg)}} \label{ham_dec}

Decodes the Hamming code using the transpose of the parity-checker
matrix, $\boldsymbol{H}^\intercal$, and passes back the decoded data
through the pointer \texttt{msg}. The decoded data bits should
constitute the lower nibble. If no correction is needed when decoding
\texttt{code}, return \texttt{HAM\_OK}. If one error is detected and
corrected, return \texttt{HAM\_CORRECT}. In the event that more than one
error is detected, return \texttt{HAM\_ERR} to indicate an uncorrectable
error without modifying \texttt{msg}.

To avoiding comparing the error syndrome with each row in
$\boldsymbol{H}^\intercal$ until a match is found (or no match is found
and the error cannot be corrected), we can refer to a lookup table. A
lookup table is an array that contains precomputed information that is
referred to often. By constructing a lookup table, we can avoid
performing the same computation many times at the expense of storing the
table in memory (in this case, storing $2^4$ or 16 bytes is
\emph{negligible}). The index to the table will be the error syndrome
and the value is the bit(start counting at 0) that needs to be flipped
if the error can be corrected. Thus, if $\Vec{e} = \begin{pmatrix}1 & 0
& 1 & 1 \end{pmatrix}$ or $1101_2$ ($13_{10}$), then
$\texttt{table[13]} = 1$. If the error cannot be corrected then
$\texttt{table[$\Vec{e}$]\;= HAM\_ERR}$.

\medskip
\begin{prelab}{Pre-lab Questions}
    \begin{enumerate}
        \item Complete the rest of the look-up table shown below.\\
        \begin{tabular}{ c|c }
            0 & 0 \\
            1 & 4 \\
            $\cdots$ & $\cdots$ \\
            15 & HAM\_ERR
        \end{tabular}
        \item Decode the following codes. If it contains an error, show
          and explain how to correct it. Remember, it is possible for a
          code to be uncorrectable.
        \begin{enumerate}
            \item $1110$ $0011_2$
            \item $1101$ $1000_2$
        \end{enumerate}
    \end{enumerate}
\end{prelab}

\section{Your Task}

You will be creating two programs for this assignment: an \emph{encoder}
and a \emph{decoder}. The encoder will generate Hamming codes given
input data and the decoder will decode the generated Hamming codes. Both
will operate with the Hamming(8, 4) systematic code. The generator
matrix $\boldsymbol{G}$ and the transpose of the parity-checker matrix,
$\boldsymbol{H}^\intercal$, will be represented with bit matrices.

They must read from \texttt{stdin} by default or a file and write to
\texttt{stdout} by default \emph{or a file}. Thus, they should support
command-line arguments, \texttt{-i} and \texttt{-o}, to specify input
and/or output files.  The decoder will also print statistics such as
total bytes processed, uncorrected errors, corrected errors, and the
error rate to \texttt{stderr}.

\subsection{Encoder Command-line Options}

Your encoder program, named \texttt{encode}, must support any
combination of the following command-line options.

\begin{itemize}
  \item \textbf{\texttt{-h}}\,: Prints out a help message describing the purpose
    of the program and the command-line options it accepts, exiting the
    program afterwards. Refer to the reference program in the resources
    repo for an idea of what to print.
  \item \textbf{\texttt{-i infile}}\,: Specify the input file path
    containing data to encode into Hamming codes. The default input
    should be set as \texttt{stdin}.
  \item \textbf{\texttt{-o outfile}}\,: Specify the output file path to
    write the encoded data (the Hamming codes) to. If not specified, the
    default output should be set as \texttt{stdout}.
\end{itemize}

\subsection{Decoder Command-line Options}

Your decoder program, named \texttt{decode}, will need to support the
same command-line options that your encoder does, plus an option to
trigger the verbose printing of statistics.

\begin{itemize}
  \item \textbf{\texttt{-h}}\,: Prints out a help message describing the purpose
    of the program and the command-line options it accepts, exiting the
    program afterwards. Refer to the reference program in the resources
    repo for an idea of what to print.
  \item \textbf{\texttt{-i infile}}\,: Specify the input file path
    containing Hamming codes to decode. The default input
    should be set as \texttt{stdin}.
  \item \textbf{\texttt{-o outfile}}\,: Specify the output file path to
    write the decoded Hamming codes to. If not specified, the default
    output should be set as \texttt{stdout}.
  \item \textbf{\texttt{-v}}\,: Prints statistics of the decoding
    process to \texttt{stderr}. The statistics to print are the total
    bytes processed, uncorrected errors, corrected errors, and the error
    rate. The error rate is defined as $(\text{uncorrected errors} /
    \text{total bytes processed})$, the ratio of uncorrected errors to
    total bytes processed.
\end{itemize}

\subsection{Example Program Output}

\begin{figure}[ht]
\begin{centering}
  \begin{lstlisting}[style=bashstyle]
 $ ./encode -i frankenstein.txt | ./decode | diff - frankenstein.txt
 Total bytes processed: 902092
 Uncorrected errors: 0
 Corrected errors: 0
 Error rate: 0.000000\end{lstlisting}
  \caption{Example usage and decoder statistics. Here \texttt{diff} uses the dash
  to represent \texttt{stdin}. }
\end{centering}
\end{figure}

We will be providing source code for a program (\texttt{error.c}) that
will inject errors (noise) into your Hamming codes. Note that not all
errors will be correctable. The rate at which the program injects errors
is a command-line argument and is specified by \texttt{-e rate} (default
is 0.01 or 1\%) and must be between [0.0, 1.0]. The seed is specified
with \texttt{-s seed} and it must be a positive integer.

\begin{figure}[h]
\begin{centering}
  \begin{lstlisting}[style=bashstyle]
 $ ./encode < frankenstein.txt | ./error -e 0.002 -s 2021 | ./decode > /dev/null
 Total bytes processed: 902092
 Uncorrected errors: 89
 Corrected errors: 14267
 Error rate: 0.000099\end{lstlisting}
  \caption{Example usage with added noise.}
\end{centering}
\end{figure}

\subsection{\texttt{FILE *}}

There is a lot of complicated logic at the hardware level to read and
write actual bits of information. This is out of scope for a person
wanting to write a software program. To make things simpler, the
\texttt{C} standard library provides functions that abstracts much of
this complicated logic. All of these functions are defined in
\texttt{stdio.h}, and utilize the \texttt{FILE} object. The
\texttt{FILE} object contains all the information about the file such as
its name, file position indicator, file access mode (read/write), and
other details necessary for I/O operations. Some functions of note for
this assignment are \texttt{fopen()}, \texttt{fclose()},
\texttt{fgetc()}, and \texttt{fputc()}. Read the \texttt{man} pages for
these functions to see what they do.

\subsection{File I/O and Permissions}

In the last assignment, graphs were read from either \texttt{stdin} or a
file, and the final paths was printed to either \texttt{stdout} or a
file. Thus, you should be familiar with file I/O. However, we must pay
special attention to the input and output file's permissions when
generating or decoding Hamming codes. For example, imagine you have a
sensitive file on your \textsc{UNIX} machine that only you can read or
write. Everyone else will receive an error if they attempt to read or
write to it. The file containing the generated Hamming codes should also
only be read or written by you, \textit{i.e.}, it should inherit the
file permissions of the file for which it is generating Hamming codes.
This will prevent another user with prying eyes to decode the file with
Hamming codes and learn the contents of your sensitive file (if it is
\emph{that} sensitive, it should be encrypted in the first place).

\subsubsection{\textsc{UNIX} File Permissions}

Every file in \textsc{UNIX} has a set of access modes for the owner,
group, and everyone else. The \texttt{ls -l} command will lists all the
files in a directory and also displays its permissions.  In figure
\ref{batPerm}, the owner of the file, \texttt{bat-notes}, is batman, and
the group is \texttt{dc}. Users in \textsc{UNIX} can be members of a
group such as the \texttt{sudo} group. The \texttt{groups} command
lists the groups the current user is in. Groups allow users to share
files within their groups while controlling the extent of their access.

\begin{figure}[h]
\begin{centering}
  \begin{lstlisting}[style=bashstyle]
 $ ls -l
   -rw-rw-r-- 1 batman dc 102 Jan 21 12:33 bat-notes\end{lstlisting}
  \caption{Output of \texttt{ls -l}}
  \label{batPerm}
\end{centering}
\end{figure}

The file's permissions are presented in the same order from left to
right after the first dash: owner, group, and others.  Read permission
is represented by \texttt{r}, write permission, the ability to modify,
create, or delete a file, is represented by \texttt{w}, and execute
permission for programs and shell scripts are represented by \texttt{x}.
Thus, in figure \ref{batPerm}, batman has the following permissions on
his file, bat-notes: \texttt{rw-rw-r--}. The owner, \texttt{batman}, has
read and write permissions but cannot execute the file. The group,
\texttt{dc}, has read and write permissions but cannot execute the file.
Everyone else can only read the file.

File permissions are not set in stone and can be modified by
\texttt{chmod} (both a command and a system call).  For example, if
batman wants to remove read access to everyone else other than himself
and members of the dc group, he can execute the command \texttt{chmod
o-r bat-notes} to remove read permissions from others. If he changes his
mind later, he can use the command \texttt{chmod o+r bat-notes} to add
read permissions to others.

\begin{figure}[h]
\begin{centering}
  \begin{lstlisting}[style=bashstyle]
 $ chmod o-r,g-w bat-notes
 $ ls -l
   -rw-r---- 1 batman dc 102 Jan 21 12:33 bat-notes\end{lstlisting}
  \caption{Removing other's read permission and the group's write permission}
  \label{batPerm1}
\end{centering}
\end{figure}

As mentioned earlier, \texttt{chmod} is also a system call and can be
called from a \textbf{C} program.  For this assignment, all output files
should inherit the file permissions of the input file. \texttt{fstat()}
should be used to retrieve the input file's permissions and
\texttt{fchmod()} to set the output file's permissions to match the
input's. Refer to \texttt{man chmod, man fchmod, man fstat} for more
information.  Note: both of these functions expect a \emph{file
descriptor} that is returned by low-level IO (\texttt{open()}), so you
will need to use \texttt{fileno()} to get the file descriptor of an open
stream.

\begin{codelisting}{Using \texttt{fstat()} and \texttt{fchmod()}}
 // Opening files to read from and write to.
 FILE *infile = fopen("bat-notes", "rb");
 FILE *outfile = fopen("bat-notes-encoded", "wb");

 // Getting and setting file permissions
 struct stat statbuf;
 fstat(fileno(infile), &statbuf);
 fchmod(fileno(outfile), statbuf.st_mode);
\end{codelisting}

Another way to represent file permissions is with three \emph{octal}
digits, one for owner, group, and others.  The octable table is shown in
table \ref{octal}. In fact, those with a keen eye will notice that file
permissions are in fact a set. Thus, if a file has its permissions set
to $755_8$ then the owner has read, write, and execute permissions and
the group and others only have read, execute permissions.

\begin{table}[h]
\centering
\begin{tabular}{ |c|c|c| }
 \hline
  Ref. & Octal & Binary \\
  \hline
  \texttt{-{}-x} & 1 & 001 \\
  \texttt{-w-} & 2 & 010 \\
  \texttt{-wx} & 3 & 011 \\
  \texttt{r-{}-} & 4 & 100 \\
  \texttt{r-x} & 5 & 101 \\
  \texttt{rw-} & 6 & 110 \\
  \texttt{rwx} & 7 & 111 \\
 \hline

 \hline
\end{tabular}
\caption{Octal and binary representation of file permissions.}
\label{octal}
\end{table}

\section{Packing and Unpacking Nibbles}

Here are some functions to aid you in getting either the low or high
nibble in a byte, as well as packing them together back into a byte.

\begin{codelisting}{Helper functions}
 // Returns the lower nibble of val
 uint8_t lower_nibble(uint8_t val) {
     return val & 0xF;
 }

 // Returns the upper nibble of val
 uint8_t upper_nibble(uint8_t val) {
     return val >> 4;
 }

 // Packs two nibbles into a byte
 uint8_t pack_byte(uint8_t upper, uint8_t lower) {
     return (upper << 4) | (lower & 0xF);
 }
\end{codelisting}

\section{Encoder Program Specifics}

In \texttt{main()} do the following:
\begin{enumerate}
    \item Parse the command-line options with \texttt{getopt()} and open
      any input and/or output files using \texttt{fopen()} and correct
      file permissions.

    \item Initialize the generator matrix $\boldsymbol{G}$ using \texttt{bm\_create()}.

    \item Read a byte from the specified file stream or \texttt{stdin}
      with \texttt{fgetc()}. \label{gen1}

    \item Generate the Hamming(8,4) codes for both the upper and lower
      nibble with \texttt{ham\_encode()} and write to the specified
      \label{gen2} file or \texttt{stdout} with \texttt{fputc()}. The
      Hamming code for the lower nibble should be written first followed
      by the code for the upper nibble. \textcolor{red}{Note: You'll
      notice this operation becomes repetitive with larger files. If
    only there was a lookup table of Hamming codes to refer to.}

    \item Repeat steps \ref{gen1} -- \ref{gen2} until all data has been
      read from the file or \texttt{stdout}.

    \item Close both the input and output files with \texttt{fclose()}
      and make sure that any memory allocated is freed.
\end{enumerate}

\section{Decoder Program Specifics}
  \begin{enumerate}
    \item Parse the command-line options with \texttt{getopt()} and open
      any input and/or output files using \texttt{fopen()} and correct
      file permissions.

    \item Initialize the transpose of the parity-checker matrix,
      $\boldsymbol{H}^\intercal$, using \texttt{bm\_create()}.

    \item Read \emph{two} bytes from the specified file stream or
      \texttt{stdin} with \texttt{fgetc()}. \label{dec1}
      \textcolor{red}{Note: The first byte read is the Hamming code for
      the lower nibble, and the second is the upper nibble.} \item For
      each byte pair read, decode the Hamming(8,4) codes for both with
      \texttt{ham\_decode()} to recover the original upper and lower
      nibbles of the message. Then, reconstruct the original byte.
      \textcolor{red}{Note: The return code from \texttt{ham\_decode()}
        should be used for statistics. Your program should count the
      number of bytes proccessed, Hamming codes that required
    correction, and Hamming codes that could not be corrected.}

    \item Write the reconstructed byte with \texttt{fputc()}.
      \textcolor{red}{Note: You should still write the byte even for
      Hamming codes that could not be corrected, leaving the data bits
    in the byte as is.} \label{dec2}

    \item Repeat steps \ref{dec1} -- \ref{dec2} until all data has been
      read from the file or \texttt{stdout}.

    \item Print the following statistics to \texttt{stderr} with \texttt{fprintf()}:

    \begin{itemize}
        \item Total bytes processed: The number of bytes read by the
          decoder.
        \item Uncorrected errors: The number of Hamming codes that could
          not be corrected.
        \item Corrected errors: The number of Hamming codes that
          experienced an error that was recoverable.
        \item Error rate: The rate of uncorrected errors for a given
          input. Refer to \S 6.2 for the calculation.
    \end{itemize}
    \item Close both the input and output files with \texttt{fclose()}
      and make sure that any memory allocated is freed.
\end{enumerate}

Remember: In \emph{both} the encoder and decoder, if an input and output
file are specified, the output file should have the same file
permissions as the input file. You can use \texttt{fstat()} to retrieve
an open file's permissions, \texttt{fchmod()} to change the permissions
of the output file to match the input's. Since both of these functions
expect a \emph{file descriptor} that is returned by low-level IO
(\texttt{open()}), you will need to use \texttt{fileno()} to get the
file descriptor of an open stream.

\section{Deliverables}

\noindent You will need to turn in:
\begin{enumerate}
  \item \texttt{encode.c}: This file will contain your implementation of the
    Hamming Code encoder.

  \item \texttt{decode.c}: This file will contain your implementation of the
    Hamming Code decoder.

  \item \texttt{error.c}: This file will be provided in the resources
    repo, but should be included in your repo as well.
    \textcolor{red}{You \emph{must not} modify this file.}

  \item \texttt{entropy.c}: This file will be provided in the resources
    repo, but should be included in your repo as well.
    \textcolor{red}{You \emph{must not} modify this file.}

  \item \texttt{bv.h}: This fill will contain the bit vector ADT interface. This file
  will be provided. \textcolor{red}{You \emph{may not} modify this file.}

  \item \texttt{bv.c}: This file will contain your implementation of the bit vector ADT.
    You \emph{must} define the bit vector \texttt{struct} in this file.

  \item \texttt{bm.h}: This file will contain the bit matrix ADT interface. This file
  will be provided. \textcolor{red}{You \emph{may not} modify this file.}

  \item \texttt{bm.c}: This file will contain your implementation of the
    bit matrix ADT. You \emph{must} define the bit matrix
    \texttt{struct} in this file.

  \item \texttt{hamming.h}: This file will contain the interface of the
    Hamming Code module. This file will be provided. \textcolor{red}{You
    \emph{may not} modify this file.}

  \item \texttt{hamming.c}: This file will contain your implementation of the
    Hamming Code module.

  \item \texttt{Makefile}: This is a file that will allow the grader to
    type \texttt{make} to compile your programs.

    \begin{itemize}
      \item \texttt{CC=clang} must be specified.
      \item \texttt{CFLAGS=-Wall -Wextra -Werror -Wpedantic}
        must be included.
      \item \texttt{make} should build the encoder, the decoder, the
        supplied error-injection program, \emph{and} the supplied
        entropy-measure program, as should \texttt{make all}.
      \item \texttt{make encode} should build \emph{just} the encoder.
      \item \texttt{make decode} should build \emph{just} the decoder.
      \item \texttt{make error} should build \emph{just} the supplied
        error-injection program.
      \item \texttt{make entropy} should build \emph{just} the supplied
        entropy-measure program.
      \item \texttt{make clean} must remove all files that are compiler
        generated.
      \item \texttt{make format} should format all your source code,
        including the header files.
    \end{itemize}

  \item Your code must pass \texttt{scan-build} \emph{cleanly}. If there
    are any bugs or errors that are false positives, document them and
    explain why they are false positives in your \texttt{README.md}.

  \item \texttt{README.md}: This must be in \emph{Markdown}. This must describe
    how to build and run your program.

  \item \texttt{DESIGN.pdf}: This \emph{must} be a PDF\@. The design document
    should answer the pre-lab questions, describe the purpose of your program,
    and communicate its overall design with enough detail such that a sufficiently
    knowledgeable programmer would be able to replicate your implementation.
    \textcolor{red}{This does not mean copying your entire program in verbatim.}
    You should instead describe how your program works with supporting pseudocode.
    \textcolor{red}{\textbf{C} code is \textbf{not} considered pseudocode.}

  \item \texttt{WRITEUP.pdf}: This document \emph{must} be a PDF. You
    will need to use the \texttt{entropy.c} program that will be
    supplied to you. The program reads data from \texttt{stdin} and
    prints out the amount of entropy in the read data. Entropy, as
    defined by Claude Shannon, quantifies the amount of information that
    is produced by some process. This will be discussed further in
    lecture. Your writeup must include the following:
    \begin{itemize}
      \item Graphs showing the amount of entropy of data before and
        after encoding it.
      \item Analysis of the graphs you produce.
    \end{itemize}
\end{enumerate}


\section{Submission}

To submit your assignment through \texttt{git}, refer to the steps shown in
\texttt{asgn0} Remember: \emph{add, commit,} and \emph{push}!
\textcolor{red}{Your assignment is turned in \emph{only} after you have pushed
and submitted the commit ID to Canvas. Your design document is turned in \emph{only}
after you have pushed and submitted the commit ID to Canvas.
If you forget to push, you have not turned in your assignment and you will get a
\emph{zero}. ``I forgot to push'' is not a valid excuse. It is \emph{highly}
recommended to commit and push your changes \emph{often}.}


\section{Supplemental Readings}

\epigraph{\emph{The more that you read, the more things you will know. The
more that you learn, the more places you'll go.}}{---Dr.\ Seuss}

\begin{itemize}
  \item \textit{The C Programming Language} by Kernighan \& Ritchie
  \begin{itemize}
    \item Chapter 2 \S 2.9
    \item Chapter 5 \S 5.7
    \item Chapter 7
  \end{itemize}
\end{itemize}

\begin{center}
  \includegraphics[width=0.35\textwidth]{../../images/monkey-chainsaw.jpg} \\
  \vspace{10pt}
  % Vigenere cipher. Key: hamming
  \texttt{Wrasznstizs qa I ps xusr mpvuzo n svnwqg n ioauzanc.}
\end{center}

\end{document}
